\section{System: \blackrim Main-Thread and Dispatch Architecture}
\label{sec:system}

\subsection{Main-thread orchestrator}

\blackrim's orchestrator, Gestalt, runs on the main conversation
thread and processes operator turns in sequence.  In the baseline
configuration, Gestalt runs on \texttt{claude-opus-4-7} for every
turn.  It decomposes operator requests into subtasks, selects the
appropriate citizen or worker agent, composes a brief, and dispatches
via the \texttt{Agent} tool.  Dispatches spawn into isolated
git worktrees (one per agent invocation) so that parallel agents do
not clobber each other's file edits.

The main thread is architecturally distinct from the dispatch path in
two ways relevant to this paper.  First, the main thread carries the
full static prefix (\claudemd instructions, citizen definitions,
project context) on every turn, making per-turn context lengths
consistently large (typically 4\,000–20\,000 tokens).  Second, the
main thread has no role taxonomy: the same Gestalt instance handles
clarification questions, architectural decisions, code review,
orchestration planning, and status summaries within a single session.
This heterogeneity is precisely what makes per-turn routing
non-trivial.

\subsection{Crew architecture}

Below Gestalt, eight \emph{citizens} (opus-tier domain leads) own
vertical domains: engineering, frontend, infrastructure, security,
product, data, ML, and knowledge.  Seven \emph{exoselves} (sonnet or
haiku) are functional workers — builder, reviewer, tester, researcher,
writer, judge, and supervisor — available to any citizen.  The
assignment of citizens and workers to tiers is fixed by domain
expertise: architect decisions stay on opus; file search and data
gathering route to haiku; code generation routes to sonnet.  This
dispatch-side routing is handled by the subagent advisor described
next; it is not the subject of this paper.

\subsection{Existing dispatch-side routing: the 16-rule cascade}

The companion paper~\cite{moa1landscape} describes the subagent
advisor (\emph{CC-TS}) that routes dispatched agents.
\texttt{internal/dispatch/model.go} implements a 16-rule
role-\(\times\)-signal cascade: rules fire in priority order, the
first match wins, and the default is sonnet.  Selected rules include:
architect role always receives opus; read-only research receives haiku;
most edit and planning work receives sonnet; novel design or
load-bearing decisions override to opus via frontmatter flag.  This
cascade is the existing baseline; it explains the 0.6\,\% dispatch
share of total session cost (§\ref{sec:eval}).

This paper's contribution is the \emph{main-thread analog} of that
cascade: a per-turn router that operates on the operator–assistant
conversation stream rather than on a structured (role, signal) pair.
Unlike the dispatch cascade, the main-thread router cannot rely on an
explicit role label; it must classify from turn content alone.

\subsection{Telemetry instrumentation}

\blackrim records per-call telemetry to
\texttt{.beads/telemetry/invocations.jsonl}.  Each record contains:
model identifier, input and output token counts, the
\texttt{cache\_creation\_input\_tokens} and
\texttt{cache\_read\_input\_tokens} fields returned by the Anthropic
API, source tag (\texttt{subagent\_stop}, \texttt{dispatch},
\texttt{gt-cache-warm}, or \texttt{estimated-stub}), and a UTC
timestamp.  This log is the substrate for §\ref{sec:eval}'s
measurements.  The 2026-05-09 through 2026-05-15 window analysed
below contains 744 spawn records, of which 81 are real-API calls with
non-zero cache fields and 663 are estimated stubs generated during
offline eval runs.

\subsection{Integration points for the new primitives}

The two primitives each intercept a distinct step in the session lifecycle:
the router acts before the first LLM call for each turn, and the plan cache
acts before the orchestration-plan composition step for each dispatch.

\paragraph{Per-turn router.}
The router sits between the operator turn arriving and Gestalt's
first LLM call for that turn.  In production, it would be implemented
as middleware in the Agent SDK before the \texttt{client.messages.create}
call.  In the current Claude Code environment, per-turn tier selection
is blocked by the platform: Claude Code exposes only session-level
model selection, not call-level routing.  §\ref{sec:discussion}
analyses three workaround paths.  For the evaluation in
§\ref{sec:eval}, routing decisions are replayed offline against the
hand-labelled exemplar set.

\paragraph{Plan cache.}
The plan cache sits between the \texttt{gt orchestrate plan} call
and the \texttt{gt context budget} briefing-composition step.
On a cache hit, the briefing-composition LLM call is bypassed entirely;
the cached plan output is injected directly into the dispatch path.
This short-circuit is why plan-cache savings are orthogonal to
prefix-cache savings: the prefix cache discounts the static-prefix
portion of LLM calls that \emph{do} fire; the plan cache prevents
certain LLM calls from firing at all.
